\documentclass[10pt,a4paper]{amsart}
\usepackage[margin=19mm]{geometry}
\usepackage{amsmath,amssymb,mathtools}
\usepackage[scr=rsfs,cal=euler]{mathalfa}
\usepackage{xfrac}
\usepackage[unicode,hidelinks,bookmarks=false]{hyperref}
\newcommand{\ie}{i.e.\ }
\newcommand{\cf}{cf.\ }
\newcommand{\resp}{respectively\ }
\newcommand{\Subsection}{\subsection}
\providecommand{\nobreakdash}{\nobreak\hbox{-}}
\setlength{\emergencystretch}{2em}
\multlinegap0pt
\usepackage{amssymb}
\newtheorem{theorem}{Theorem}
\newcommand{\Z}{\mathbb{Z}}
\title{A proper Euler magic matrix of order $6$}
\author{Sanjit Singh Mehat}
\date{Source: arXiv:2608.15318v1 (CC BY 4.0)}
\begin{document}
\maketitle

\noindent\emph{Source and licence.} A proper Euler magic matrix of order $6$. Version arXiv:2608.15318v1 (CC BY 4.0), distributed by arXiv under the Creative Commons Attribution 4.0 International licence, http://creativecommons.org/licenses/by/4.0/, as recorded in the arXiv licence metadata for this exact version and shown on the abstract page https://arxiv.org/abs/2608.15318v1. Full licence notice: \texttt{licenses/arxiv-2608.15318-CC-BY-4.0.txt}.\par\smallskip

\noindent\textit{Editorial scope.} Counted unit: the article's theorem thm:main (source0.tex lines 69--83). The article's complete original proof of it is delivered with it (source0.tex lines 85--109, one \texttt{proof} environment of 25 lines). No mathematics is written, repaired or re-derived: the statement and the proof are the article's own lines, in the article's own notation, and no retained line is substituted or re-flowed.  Retained uncounted as the source-local context the proof needs: lines 44--48.  Nothing was omitted from any retained span, so the package carries no omission record.  No retained line contains a citation, so the wrapper carries no bibliography and no bibliography entry.  No retained line contains a figure environment, an image-inclusion command or drawing code, so no image is delivered and the package declares no asset dependency.  Editorial formatting only, in addition: an AMS \texttt{amsart} preamble that restates the article's own declarations and macros used by the retained lines, namely the environments \texttt{theorem} on separate counters, together with the standard packages those lines need.  The retained environments print in this document as Theorem 1 in order of appearance, whereas the article prints the same items with the numbering of its own full document.  The complete native source of the article as distributed in the arXiv e-print for this exact version is bundled unchanged as \texttt{sources/207731/source0.tex}.
\begin{align}
M\cdot M^{t} &= \gamma\cdot I_n, \label{eq:orth}\\
\textstyle\sum_{i=1}^{n} m_{i,i}^{2} &= \gamma, \label{eq:diag}\\
\textstyle\sum_{i=1}^{n} m_{i,n+1-i}^{2} &= \gamma. \label{eq:anti}
\end{align}

\begin{theorem}\label{thm:main}
A proper Euler magic matrix of order $6$ exists. Explicitly, the matrix
\[
M \;=\;
\begin{pmatrix}
-55 & -52 & -89 & -53 & -21 &  40\\
-30 & -83 &  49 &  -2 &  91 &   5\\
 69 & -71 & -42 &  81 & -18 &   7\\
 90 &   9 & -37 & -74 &  57 & -15\\
-33 &  -8 & -41 &  13 &  11 &-124\\
 25 & -61 &  58 & -59 & -78 & -35
\end{pmatrix}
\]
is an Euler magic matrix with $\gamma=18500$, and it is proper.
\end{theorem}

\begin{proof}
All three conditions are finite identities in $\Z$ and are verified directly.
For \eqref{eq:orth}, each row has squared norm
\[
55^2+52^2+89^2+53^2+21^2+40^2=18500,
\]
and similarly for the remaining five rows, while the fifteen distinct pairs of
rows are orthogonal. For \eqref{eq:diag},
\[
55^2+83^2+42^2+74^2+11^2+35^2 = 3025+6889+1764+5476+121+1225 = 18500 ,
\]
and for \eqref{eq:anti},
\[
40^2+91^2+81^2+37^2+8^2+25^2 = 1600+8281+6561+1369+64+625 = 18500 .
\]
Finally, the $36$ absolute values
\[
\begin{aligned}
\{\,&2,5,7,8,9,11,13,15,18,21,25,30,33,35,37,40,41,42,\\
 &49,52,53,55,57,58,59,61,69,71,74,78,81,83,89,90,91,124\,\}
\end{aligned}
\]
are pairwise distinct, and over $\Z$ one has $x^2=y^2$ if and only if
$|x|=|y|$; hence the $36$ squares are pairwise distinct and $M$ is proper.
\end{proof}

\end{document}
